%! Adding \& Subtracting Rational Expressions — Unit 5, Lesson 5.3 — Group Activity (Answer Key)
\documentclass[10pt]{article}
\usepackage{algebra2-article}
\usepackage{algebra2-key}   % pulls in -boxes; do NOT also load -boxes
\newcommand{\TallMath}[1]{$\displaystyle #1\rule[-1.4em]{0pt}{3.2em}$}

\begin{document}

\pageheader{Unit 5, Lesson 5.3}{Group Activity --- Answer Key}
\namepartnerperiod

\vspace{0.08in}

\begin{headlinebox}{forestbg}
\small\textbf{Build the Denominator, Guard the Sign.} Two habits carry everything today. Factor every
denominator \emph{before} you decide what the LCD is --- it is built out of factors, never out of
products. And every time you subtract, put the second numerator in parentheses \emph{first}. As always,
an answer without its restrictions is half an answer.
\end{headlinebox}

\vspace{0.06in}

\begin{tcolorbox}[colback=white, colframe=black!40, title=\textbf{Tier R --- Remediate},
    fonttitle=\bfseries, arc=1.5mm, left=3mm, right=3mm, top=2mm, bottom=2mm, breakable]
Combine and simplify. State every restriction.

\vspace{4pt}
\begin{enumerate}[leftmargin=1.9em, itemsep=9pt]
  \item $\dfrac{2x}{x-6}+\dfrac{5}{x-6}$ \; (denominators already match) \\[3pt]
        Simplest form: \ans{$\frac{2x+5}{x-6}$} \quad $x\neq$ \ans{6}
  \item $\dfrac{4x+1}{x+2}-\dfrac{x-8}{x+2}$ \; Write the numerator with parentheses first:
        $(4x+1)-({\color{keyred}\mathbf{x-8}})$ \\[3pt]
        Numerator, simplified and factored: \ans{$3x+9=3(x+3)$} \quad Simplest form: \ans{$\frac{3(x+3)}{x+2}$}
        \quad $x\neq$ \ans{$-2$}
  \item $\dfrac{3}{x}+\dfrac{5}{2x}$ \\[3pt]
        LCD: \ans{$2x$} \quad Building factor for the first fraction: \ans{$\frac22$} \quad
        Answer: \ans{$\frac{11}{2x}$}, \; $x\neq$ \ans{0}
  \item $\dfrac{1}{x-1}+\dfrac{2}{x+3}$ \\[3pt]
        LCD: \ans{$(x-1)(x+3)$} \quad Answer: \ans{$\frac{3x+1}{(x-1)(x+3)}$} \quad $x\neq$ \ans{$1,\ -3$}
  \item In problem 2, do it the wrong way on purpose: drop the parentheses and write
        $4x+1-x-8$. \\[3pt]
        Wrong numerator: \ans{$3x-7$} \quad Now test both at $x=0$: the original expression equals
        \ans{$\frac92$}, your correct answer equals \ans{$\frac92$}, and the wrong one equals
        \ans{$-\frac72$}.
\end{enumerate}
\end{tcolorbox}

\begin{tcolorbox}[colback=white, colframe=black!40, title=\textbf{Tier A --- Approaching Proficiency},
    fonttitle=\bfseries, arc=1.5mm, left=3mm, right=3mm, top=2mm, bottom=2mm, breakable]
\textbf{Part 1 --- The full table.} Fill in every cell. One row hides a pair of \emph{opposite}
binomials, and one row simplifies further than you expect.
\begin{center}
\renewcommand{\arraystretch}{1.7}
\begin{tabularx}{\linewidth}{@{}l X X X@{}}
\hline
\textbf{Expression} & \textbf{Denominators, factored} & \textbf{LCD \& restrictions} & \textbf{Simplest form} \\
\hline
\TallMath{\frac{7}{4x^2}-\frac{3}{10x}} & \ans{$2^2x^2$ and $2\cdot5x$} & \ans{$20x^2$; $x\neq0$} & \ans{$\frac{35-6x}{20x^2}$} \\
\TallMath{\frac{x}{x^2-4}+\frac{1}{x+2}} & \ans{$(x-2)(x+2)$ and $(x+2)$} & \ans{$(x-2)(x+2)$; $x\neq2,-2$} & \ans{$\frac{2(x-1)}{(x-2)(x+2)}$} \\
\TallMath{\frac{2}{x-4}-\frac{16}{x^2-16}} & \ans{$(x-4)$ and $(x-4)(x+4)$} & \ans{$(x-4)(x+4)$; $x\neq4,-4$} & \ans{$\frac{2}{x+4}$} \\
\TallMath{\frac{5}{x-2}+\frac{3}{2-x}} & \ans{$(x-2)$ and $-(x-2)$} & \ans{$x-2$; $x\neq2$} & \ans{$\frac{2}{x-2}$} \\
\hline
\end{tabularx}
\end{center}
{\small Row 3 is the one that ``simplifies further than you expect'': the numerator $2x-8$ factors as
$2(x-4)$ and divides out, so $x=4$ becomes a \textbf{hole} while $x=-4$ stays a \textbf{wall}. Row 4 is
the opposite-binomial row --- its LCD is $x-2$, \emph{not} $(x-2)(2-x)$.}

\vspace{4pt}
\textbf{Part 2 --- Error analysis.} A student turns in this work:
\begin{center}
\fbox{\parbox{0.66\linewidth}{\centering \vspace{2pt}
$\dfrac{1}{x} + \dfrac{1}{3} \;=\; \dfrac{1+1}{x+3} \;=\; \dfrac{2}{x+3}$ \\[3pt]
{\small ``I added the tops and I added the bottoms.''}\vspace{2pt}}}
\end{center}
\begin{enumerate}[label=(\alph*), leftmargin=1.9em, itemsep=6pt]
  \item Disprove it with a number. At $x=1$: the original equals \ans{$\frac43$} and the student's
        answer equals \ans{$\frac12$}.
  \item Do it correctly: LCD $=$ \ans{$3x$}, so $\dfrac1x+\dfrac13=$ \ans{$\frac{x+3}{3x}$}, \;
        $x\neq$ \ans{0}.
  \item Multiplying rational expressions \emph{does} let you work straight across the tops and the
        bottoms. Adding does not. In one sentence, what is different about adding?
        \ansline{A denominator says how big each piece is. Multiplying makes new pieces, so the denominators}
        \ansline{legitimately multiply; adding only counts pieces you already have, so they must be the same size first.}
\end{enumerate}
\end{tcolorbox}

\begin{tcolorbox}[colback=white, colframe=black!40, title=\textbf{Tier E --- Extension},
    fonttitle=\bfseries, arc=1.5mm, left=3mm, right=3mm, top=2mm, bottom=2mm, breakable]
\textbf{Part 1 --- One complex fraction, two roads.} Simplify
\[
  \frac{\;\dfrac{1}{x}-\dfrac{1}{4}\;}{\;\dfrac{1}{x^2}-\dfrac{1}{16}\;}
\]
\begin{enumerate}[label=(\alph*), leftmargin=1.9em, itemsep=6pt]
  \item \textbf{Method 1 (combine, then divide).} Top: \ans{$\frac{4-x}{4x}$} \quad Bottom: \ans{$\frac{(4-x)(4+x)}{16x^2}$}
        \\[3pt]
        Keep--change--flip and simplify: \ans{$\frac{4-x}{4x}\cdot\frac{16x^2}{(4-x)(4+x)}=\frac{4x}{4+x}$}
  \item \textbf{Method 2 (multiply through by the LCD of the little fractions).} That LCD is
        \ans{$16x^2$}. \\[3pt]
        Top becomes \ans{$16x-4x^2$} ; bottom becomes \ans{$16-x^2$} ; answer: \ans{$\frac{4x}{4+x}$}
  \item Complete restriction list: $x\neq$ \ans{$0,\ 4,\ -4$}. Which value(s) came from the \emph{main
        fraction bar}? \ans{$4$ and $-4$}
  \item One method leaves \emph{every} restriction readable somewhere on the page; the other erases
        one of them. Which method erases which restriction, and what step destroyed the evidence?
        \ansline{Method 2 erases $x\neq0$: multiplying top and bottom by $16x^2$ clears every inner denominator, so}
        \ansline{no $x$ is left downstairs to warn you. Method 1 keeps $4x$ and $16x^2$ in view, and the bottom's numerator}
        \ansline{$(4-x)(4+x)$ is the divisor's numerator from Lesson 5.2 --- which is where $x\neq4,-4$ come from.}
\end{enumerate}

\vspace{5pt}
\textbf{Part 2 --- Build one backwards.} Write a \emph{difference} of two rational expressions whose
simplified value is $\dfrac{2}{x+1}$ and whose restrictions are exactly $x\neq-1$ and $x\neq4$.
\begin{center}
\ans{$\frac{2}{x-4}$} $-$ \ans{$\frac{10}{x^2-3x-4}$}
\end{center}
\begin{enumerate}[label=(\alph*), leftmargin=1.9em, itemsep=5pt]
  \item What is the LCD of your two denominators, and why did it have to contain $(x-4)$?
  \item On the graph of your expression, which of the two excluded values is a \textbf{hole} and which
        is a \textbf{wall}? How do you know without graphing anything?
\end{enumerate}
\ansline{(a) $x^2-3x-4=(x-4)(x+1)$, so the LCD is $(x-4)(x+1)$ --- it must contain $(x-4)$ because that is the}
\ansline{first denominator. Combining gives $\frac{2(x+1)-10}{(x-4)(x+1)}=\frac{2(x-4)}{(x-4)(x+1)}=\frac{2}{x+1}$. (b) $x=4$ is a \textbf{hole}, because}
\ansline{its factor divided out; $x=-1$ is a \textbf{wall}, because $(x+1)$ is still a denominator in the answer.}

\vspace{5pt}
\textbf{Part 3 --- Two crews, one mural.} Crew A can paint a mural alone in $x$ hours. Crew B is slower
and needs $x+3$ hours.
\begin{enumerate}[label=(\alph*), leftmargin=1.9em, itemsep=6pt]
  \item In one hour, Crew A finishes \ans{$\frac1x$} of the mural and Crew B finishes \ans{$\frac{1}{x+3}$} of
        it. Working together, they finish \ans{$\frac{2x+3}{x(x+3)}$} of it per hour.
  \item Time is $1$ whole mural divided by their combined rate --- a \textbf{complex fraction}.
        Simplify it: \\[3pt]
        Time together $=$ \ans{$\frac{x(x+3)}{2x+3}$} hours.
  \item If Crew A alone takes $x=6$ hours, the two crews together take \ans{$3.6$} hours. Is that
        answer smaller than both $6$ and $9$? \ans{yes} Why must it be? \ans{help never slows you down}
  \item The algebra excludes $x=0$, $x=-3$, and $x=-\tfrac32$. Which values of $x$ actually make sense
        for two paint crews? \ans{$x>0$}
\end{enumerate}
\end{tcolorbox}

\begin{teachernote}
\textbf{Tier R} is deliberately front-loaded with the sign work: item 2 is the whole subtraction rule,
and item 5 makes students produce the \emph{wrong} answer on purpose and then watch a number refuse to
match. Do not let them skip item 5 --- it is the only place the error becomes personal.

\vspace{3pt}
\textbf{Tier A}, row 3 of the table, is where the day's ideas meet: an addition problem whose answer
cancels, producing a hole from a denominator that no longer appears. Row 4 catches the group that
builds an LCD of $(x-2)(2-x)$ --- ask them what that product actually equals. Part 2 is the second
signature error; require the number at $x=1$ before any explanation, since ``you can't do that'' is not
a disproof.

\vspace{3pt}
\textbf{Tier E, Part 1(d)} is the conceptual payoff and worth a whole-class share-out: the fast method
is fast precisely because it throws away the inner denominators, which are also the evidence for one of
the restrictions. That is the same trade students met in 5.1 (cancelling) and 5.2 (flipping), now for
the third time --- name the pattern out loud. \textbf{Part 3} is A2.F.2f in miniature; part (c) is the
sanity check that makes the model believable, and part (d) separates an algebraic restriction from a
contextual domain. Groups that finish early can be asked what happens to the together-time as $x$ grows
large.
\end{teachernote}

\end{document}
